\section{Introduction}
\label{sec:introduction}

Critical transitions in dynamical systems are often preceded by measurable
changes in the observable time series. The practical question is whether
these changes can be detected early enough to support intervention. This
makes early-warning detection a problem of timing, calibration, and
false-alarm control---not only binary classification.

Classical early-warning indicators based on variance, autocorrelation, and
related statistics have been applied across ecology, climate science, and
engineering \citep{scheffer2009early,dakos2008slowing}. These indicators are
transparent and require no model fitting, but each encodes a fixed functional
form for the warning signature. When the transition mechanism changes, a
statistic that was informative for one bifurcation type may become unreliable
for another.

Learned observers offer an alternative. A Kalman-style observer maintains a
latent state updated by a prediction--correction rule and can learn the
mapping from observations to warning scores from data. Recent work has shown
that such models can extract early-warning information from simulated
systems and, in some cases, transfer across dynamics with matching
normal-form structure \citep{bury2021deep,deb2022machine}. However, a
fold bifurcation, a Hopf bifurcation, and a period-doubling transition
arise from different local mechanisms
\citep{kuznetsov2004elements,strogatz2015nonlinear} and produce different
pre-transition signatures. There is no reason to expect uniform transfer
across all of them.

This paper studies the transfer question directly. We compare two learned
Kalman observer variants---a binary-cross-entropy observer (Kalman-BCE) and
a spectral recurrent variant (Kalman-LSTM-Spec)---across synthetic fold,
Hopf, and logistic systems. The benchmark varies the training-set size from
100 to 500 trajectories and evaluates each setting across ten random seeds.
Performance is measured by three metrics: detection time, early-warning AUC,
and false-positive rate. This design separates ranking quality from
operational calibration.

The main finding is regime dependence. At 500 trajectories, the spectral
recurrent observer achieves higher early-warning AUC on fold
($0.903 \pm 0.059$ vs.\ $0.827 \pm 0.235$) and Hopf
($0.823 \pm 0.243$ vs.\ $0.617 \pm 0.177$), while both methods reach
perfect AUC on the logistic system. However, the BCE observer produces
earlier detections on fold (detection time $109.5 \pm 28.6$ vs.\
$72.7 \pm 11.0$ steps) and lower false-positive rates on fold and logistic.
The relative ranking also shifts with training-set size: LSTM-Spec benefits
more from additional data on fold and Hopf, while logistic performance
saturates by 200 trajectories for both methods. No single observer dominates
across all three systems, metrics, and data sizes.

\subsection{Contributions}
\label{sec:contributions}

We make four contributions.

\begin{enumerate}[leftmargin=*]
  \item A controlled benchmark for learned Kalman observers across fold,
        Hopf, and logistic bifurcation systems, with deterministic seed
        generation and batch selection for reproducibility.
  \item Evaluation across five training-set sizes (100--500 trajectories)
        with ten seeds per setting, providing cross-seed variance estimates
        for every reported metric.
  \item A multi-metric comparison using detection time, early-warning AUC,
        and false-positive rate, exposing trade-offs that a single aggregate
        score would hide.
  \item Evidence that observer performance is system- and metric-dependent:
        the spectral recurrent observer improves discrimination on fold and
        Hopf but does not dominate the BCE observer across all systems,
        data sizes, and metrics.
\end{enumerate}

These contributions support a bounded claim: learned observer design for
early-warning detection should be evaluated as a regime-dependent modeling
problem, not as a search for a detector that transfers uniformly across
bifurcation mechanisms.
